% GENERATED FILE. Edit canonical lessons, then run npm run generate:books.
% canonical-source-hash: fnv1a64:6ce9b00aab39d63d
% canonical-lessons: TE-C262-rotte, TE-C262-capati, TE-C262-upma, TE-C262-tiphinu, TE-C262-pesalu

\chapter{Bread, Chapati, Upma, Breakfast, Green gram}
\label{ch:te-bread-chapati-upma-breakfast-green-gram}
\hlchaptermodality{262}

\begin{chapteropening}
\textbf{By the end of this chapter:} \emph{I can say bread, a chapati, upma, breakfast and green gram.}

Complete the last lesson of chapter 262: I can say bread, a chapati, upma, breakfast and green gram.
\end{chapteropening}

\section[\texorpdfstring{\te{రొట్టె} (roṭṭe)}{రొట్టె (roṭṭe)}]{\te{రొట్టె} (roṭṭe) — bread, flatbread}

\label{lesson:TE-C262-rotte}



\begin{quote}
Before the new one: say the Telugu for sticky, then the Telugu for rough.
\end{quote}

\subsection*{You'll want to know: \te{రొట్టె}}
\textbf{\te{రొట్టె}} — \emph{roṭṭe} — \textquotedblleft{}bread, flatbread\textquotedblright{}.

\begin{grammarlens}[title={What you've built}]
One more word for this chapter.
\end{grammarlens}

\subsection*{Your turn}
\begin{itemize}
\raggedright
  \item \emph{Say it:} \emph{roṭṭe}
  \item \emph{Say it:} \emph{roṭṭe}, once more
  \item \emph{Say it:} say \emph{roṭṭe}
  \item \emph{Recall:} say the Telugu for sticky, then the Telugu for rough, then say \emph{roṭṭe} again
\end{itemize}

\begin{tcolorbox}[breakable,colback=teal!4,colframe=teal!35!black,title={Before you move on}]
What does \te{రొట్టె} mean? (\textquotedblleft{}Bread, flatbread\textquotedblright{}.) Say it once more.
\end{tcolorbox}
\section[\texorpdfstring{\te{చపాతీ} (capātī)}{చపాతీ (capātī)}]{\te{చపాతీ} (capātī) — a chapati}

\label{lesson:TE-C262-capati}



\begin{quote}
Before the new one: say the Telugu for rough, then the Telugu for bread.
\end{quote}

\subsection*{You'll want to know: \te{చపాతీ}}
\textbf{\te{చపాతీ}} — \emph{capātī} — \textquotedblleft{}a chapati\textquotedblright{}.

\begin{grammarlens}[title={What you've built}]
One more word for this chapter.
\end{grammarlens}

\subsection*{Your turn}
\begin{itemize}
\raggedright
  \item \emph{Say it:} \emph{capātī}
  \item \emph{Say it:} \emph{capātī}, once more
  \item \emph{Say it:} say \emph{capātī}
  \item \emph{Recall:} say the Telugu for rough, then the Telugu for bread, then say \emph{capātī} again
\end{itemize}

\begin{tcolorbox}[breakable,colback=teal!4,colframe=teal!35!black,title={Before you move on}]
What does \te{చపాతీ} mean? (\textquotedblleft{}A chapati\textquotedblright{}.) Say it once more.
\end{tcolorbox}
\section[\texorpdfstring{\te{ఉప్మా} (upmā)}{ఉప్మా (upmā)}]{\te{ఉప్మా} (upmā) — upma}

\label{lesson:TE-C262-upma}



\begin{quote}
Before the new one: say the Telugu for bread, then the Telugu for a chapati.
\end{quote}

\subsection*{You'll want to know: \te{ఉప్మా}}
\textbf{\te{ఉప్మా}} — \emph{upmā} — \textquotedblleft{}upma\textquotedblright{}.

\begin{grammarlens}[title={What you've built}]
One more word for this chapter.
\end{grammarlens}

\subsection*{Your turn}
\begin{itemize}
\raggedright
  \item \emph{Say it:} \emph{upmā}
  \item \emph{Say it:} \emph{upmā}, once more
  \item \emph{Say it:} say \emph{upmā}
  \item \emph{Recall:} say the Telugu for bread, then the Telugu for a chapati, then say \emph{upmā} again
\end{itemize}

\begin{tcolorbox}[breakable,colback=teal!4,colframe=teal!35!black,title={Before you move on}]
What does \te{ఉప్మా} mean? (\textquotedblleft{}Upma\textquotedblright{}.) Say it once more.
\end{tcolorbox}
\section[\texorpdfstring{\te{టిఫిను} (ṭiphinu)}{టిఫిను (ṭiphinu)}]{\te{టిఫిను} (ṭiphinu) — breakfast, a light snack}

\label{lesson:TE-C262-tiphinu}



\begin{quote}
Before the new one: say the Telugu for a chapati, then the Telugu for upma.
\end{quote}

\subsection*{You'll want to know: \te{టిఫిను}}
\textbf{\te{టిఫిను}} — \emph{ṭiphinu} — \textquotedblleft{}breakfast, a light snack\textquotedblright{}.

\begin{grammarlens}[title={What you've built}]
One more word for this chapter.
\end{grammarlens}

\subsection*{Your turn}
\begin{itemize}
\raggedright
  \item \emph{Say it:} \emph{ṭiphinu}
  \item \emph{Say it:} \emph{ṭiphinu}, once more
  \item \emph{Say it:} say \emph{ṭiphinu}
  \item \emph{Recall:} say the Telugu for a chapati, then the Telugu for upma, then say \emph{ṭiphinu} again
\end{itemize}

\begin{tcolorbox}[breakable,colback=teal!4,colframe=teal!35!black,title={Before you move on}]
What does \te{టిఫిను} mean? (\textquotedblleft{}Breakfast, a light snack\textquotedblright{}.) Say it once more.
\end{tcolorbox}
\section[\texorpdfstring{\te{పెసలు} (pesalu)}{పెసలు (pesalu)}]{\te{పెసలు} (pesalu) — green gram}

\label{lesson:TE-C262-pesalu}



\begin{quote}
Before the new one: say the Telugu for upma, then the Telugu for breakfast.
\end{quote}

\subsection*{You'll want to know: \te{పెసలు}}
\textbf{\te{పెసలు}} — \emph{pesalu} — \textquotedblleft{}green gram\textquotedblright{}.

\begin{grammarlens}[title={What you've built}]
One more word for this chapter.
\end{grammarlens}

\subsection*{Your turn}
\begin{itemize}
\raggedright
  \item \emph{Say it:} \emph{pesalu}
  \item \emph{Say it:} \emph{pesalu}, once more
  \item \emph{Say it:} say \emph{pesalu}
  \item \emph{Recall:} say the Telugu for upma, then the Telugu for breakfast, then say \emph{pesalu} again
\end{itemize}

\begin{tcolorbox}[breakable,colback=teal!4,colframe=teal!35!black,title={Before you move on}]
What does \te{పెసలు} mean? (\textquotedblleft{}Green gram\textquotedblright{}.) Say it once more.
\end{tcolorbox}
